% !TeX root = ../Lua-document-Template.tex 
% !TeX spellcheck = uk_UA

\section{Поздовжній удар вантажу по стрижню}

Розглянемо циліндричний стержень, один кінець якого ($x=0$) закріплений, а другий ($x=l$) вільний. У початковий момент часу $t=0$ вільний кінець зазнає удару вантажу маси $M$, що рухається уздовж осі стержня зі швидкістю $v$. Вивчимо поздовжні коливання стержня, що виникають під час удару по ньому.

Ми знаємо, що рівняння поздовжніх коливань однорідного стержня має вигляд
\begin{equation}
	\frac{\partial^2 u}{\partial t^2} = a^2 \frac{\partial^2 u}{\partial x^2} \quad \left(a = \sqrt{\frac{E}{\rho}}\right).
\end{equation}

Гранична умова на лівому кінці ($x=0$) буде, очевидно,
\begin{equation}
	u(0, t) = 0.
\end{equation}

Далі, рівняння руху вантажу під дією сили реакції стержня, яка дорівнює за величиною зусиллю в перетині $x=l$ стержня і спрямована в протилежну сторону, має вигляд
\begin{equation}
	M \left.\frac{\partial^2 u}{\partial t^2}\right|_{x=l} = -ES \left.\frac{\partial u}{\partial x}\right|_{x=l}.
\end{equation}

Це і буде гранична умова на кінці $x=l$. Рівнянню (3) можна надати вигляду
\begin{equation}
	ml \left.\frac{\partial^2 u}{\partial t^2}\right|_{x=l} = -a^2 \left.\frac{\partial u}{\partial x}\right|_{x=l},
\end{equation}
якщо позначити через $m = \frac{M}{\rho S}$ відношення маси вантажу, що рухається, до маси стержня.
